%==============================================================================
%  Thesis Defense -- Theoretical Foundations (Study Companion)
%  Hybrid Post-Quantum Key Exchange for WireGuard (X25519 + ML-KEM-768)
%  Mihai-Cristian Farcas
%
%  A revision companion to read side-by-side with the thesis. Each topic is a
%  self-contained briefing; cross-references point to the thesis section where
%  the concept is used.
%
%  Build:  tectonic foundations.tex
%==============================================================================
\documentclass[11pt]{article}

\usepackage[a4paper,margin=2.3cm]{geometry}
\usepackage{xcolor}
\usepackage{booktabs}
\usepackage{enumitem}
\usepackage{parskip}
\usepackage{amsmath}
\usepackage{amssymb}
\usepackage{titlesec}
\usepackage[hidelinks]{hyperref}

\definecolor{deepcyan}{HTML}{0C6E6E}
\definecolor{softgray}{HTML}{555555}
\definecolor{boxbg}{HTML}{E9F2F2}
\definecolor{rulegray}{HTML}{BBBBBB}

\hypersetup{
  pdftitle={Theoretical Foundations -- Defense Study Companion},
  pdfauthor={Mihai-Cristian Farcas},
  colorlinks=true, linkcolor=deepcyan, urlcolor=deepcyan,
}

% ---- section styling ---------------------------------------------------------
\titleformat{\section}
  {\Large\bfseries\color{deepcyan}}{\thesection}{0.6em}{}
\titleformat{\subsection}
  {\large\bfseries\color{deepcyan}}{\thesubsection}{0.6em}{}
\titlespacing*{\section}{0pt}{1.6em}{0.6em}

% ---- "maps to thesis" line ---------------------------------------------------
\newcommand{\seethesis}[1]{%
  {\small\color{softgray}\textit{In the thesis: #1}}\par\vspace{3pt}}

% ---- highlight box for likely committee questions ----------------------------
\newcommand{\keyq}[1]{%
  \par\vspace{0.4em}%
  \noindent\colorbox{boxbg}{%
    \parbox{\dimexpr\textwidth-2\fboxsep}{%
      \small\textbf{\color{deepcyan}If they probe.}\ #1}}%
  \par\vspace{0.5em}}

% ---- one-line "in one sentence" lead -----------------------------------------
\newcommand{\nutshell}[1]{\noindent\textit{In one sentence:} #1\par\vspace{4pt}}

\setlist[itemize]{leftmargin=1.4em,itemsep=2pt,topsep=2pt}
\setlist[enumerate]{leftmargin=1.6em,itemsep=2pt,topsep=2pt}
\setcounter{tocdepth}{2}

\pagestyle{plain}

\begin{document}

%--- title --------------------------------------------------------------------
\begin{center}
  {\LARGE\bfseries\color{deepcyan}Theoretical Foundations}\\[3pt]
  {\large\bfseries\color{deepcyan}A Defense Study Companion}\\[8pt]
  {\large Hybrid Post-Quantum Key Exchange for WireGuard (X25519 + ML-KEM-768)}\\[6pt]
  {\normalsize Mihai-Cristian Farca\c{s} \quad\textbar\quad Supervisor: Prof.\ Dr.\ Darabant Sergiu Adrian}
\end{center}

\vspace{0.6em}
\noindent\textbf{How to use this.} This is a briefing book, not a second thesis.
Each entry is short on purpose: enough to answer a committee question with
confidence and to follow up one level deeper if pressed. Read it alongside the
thesis --- every entry names the section(s) where the concept actually does work
in your argument. The grey \textit{``In the thesis''} line is your jump label;
the shaded \textbf{``If they probe''} boxes pre-empt the obvious follow-up.

\vspace{0.4em}
\noindent\textbf{The one-paragraph anchor.} Everything below serves a single
claim: \emph{WireGuard's key exchange relies on X25519, an elliptic-curve
Diffie--Hellman that a large quantum computer would break via Shor's algorithm;
this thesis adds ML-KEM-768 (FIPS 203) alongside X25519 in a hybrid
``concatenate-then-KDF'' construction inside BoringTun, behind a \texttt{pq}
compile flag, so that a session stays confidential as long as \emph{either}
primitive holds --- and quantifies the cost on real hardware.} Keep that
sentence in your head; each topic is a spoke off it.

\vspace{0.4em}
{\small\color{softgray}\noindent Build: \texttt{tectonic foundations.tex}. Theme matches the slide deck and speaker script.}

\vspace{0.8em}
\hrule height 1pt
\tableofcontents
\vspace{0.6em}
\hrule height 1pt

%==============================================================================
\section{Networking and VPN Foundations}
%==============================================================================
This first group locates your work on the map: \emph{where} in the stack
WireGuard sits, what it competes with, and the packet-size constraints that
post-quantum messages bump into.

%------------------------------------------------------------------------------
\subsection{The OSI model and layer 3}
\seethesis{\S2.2 (WireGuard's design); motivates why the work is transparent to applications.}
\nutshell{WireGuard operates at the network layer (layer 3), so it secures whole
IP packets and everything above them, invisibly to applications.}

The OSI reference model is a seven-layer abstraction of networking: physical (1),
data link (2), network (3), transport (4), session (5), presentation (6), and
application (7). The two reference points you need are \textbf{layer 2}, which
moves frames between directly-connected nodes (Ethernet, MAC addresses), and
\textbf{layer 3}, the \emph{network} layer, which routes \emph{packets} across
networks using IP addresses.

WireGuard is a \textbf{layer-3} tunnel: it presents a virtual network interface,
takes whole IP packets, encrypts each one, and ships it inside a UDP datagram to
the peer. Because it works on IP packets, any transport- or application-layer
protocol (TCP, QUIC, HTTP, DNS) rides over it without modification. This is the
contrast with OpenVPN's TAP mode, which can tunnel at layer 2 (Ethernet frames).
The practical consequence for your thesis: your change touches only how the
\emph{tunnel key} is established; the data-plane that carries layer-3 traffic is
untouched.

%------------------------------------------------------------------------------
\subsection{WireGuard}
\seethesis{\S2.2 and \S2.2.1 (protocol + PSK slot); the whole thesis modifies it.}
\nutshell{A deliberately minimal layer-3 VPN: one round-trip handshake built on
the Noise framework, a fixed modern cipher suite, and roughly 4{,}000 lines of
code.}

WireGuard (Jason Donenfeld, 2017; in the mainline Linux kernel since 5.6) is a
VPN built around radical simplicity. Its design choices, each of which you should
be able to name:
\begin{itemize}
  \item \textbf{Fixed cryptography, no negotiation.} Unlike TLS or IPsec there
  is no cipher-suite negotiation --- and therefore no downgrade attacks. The
  suite is Curve25519 (X25519) for key agreement, ChaCha20-Poly1305 for
  authenticated encryption, BLAKE2s for hashing and keyed MAC, and an
  HMAC-based KDF. This rigidity is exactly why adding post-quantum security
  means \emph{editing the protocol}, not flipping a negotiated option.
  \item \textbf{1-RTT Noise IK handshake.} One message each way establishes a
  session; see \S3 below.
  \item \textbf{Cryptokey routing.} Each peer is bound to a set of
  ``allowed IPs''; the public key \emph{is} the identity, which doubles as the
  routing decision.
  \item \textbf{Stateless under load / DoS resistance.} A cookie-reply
  mechanism (the \texttt{mac1}/\texttt{mac2} fields) lets a loaded responder
  defer state allocation until a client proves it can receive.
  \item \textbf{Silence is a virtue.} It runs over UDP and does not reply to
  unauthenticated packets.
\end{itemize}
BoringTun is Cloudflare's \emph{userspace} Rust implementation of this protocol
--- the codebase you extended.

\keyq{``Why is no negotiation a security feature?'' Negotiation is where
downgrade attacks live (e.g.\ TLS's history). WireGuard's fixed suite removes
that class entirely; the trade-off is that algorithm changes require a versioned
protocol change --- which is precisely the surface your \texttt{pq} flag and new
message types~5/6 add.}

%------------------------------------------------------------------------------
\subsection{IPsec}
\seethesis{Background contrast to WireGuard; RFC 8784 parallels your PSK-slot baseline.}
\nutshell{The older, standardized, kernel-level layer-3 VPN suite --- powerful
but complex, with a separate key-exchange protocol (IKEv2).}

IPsec is a suite (not one protocol) that secures IP at layer 3 using two data
formats --- ESP (Encapsulating Security Payload, confidentiality + integrity)
and AH (Authentication Header, integrity only) --- keyed by the \textbf{IKEv2}
handshake. It is mature and ubiquitous but notoriously complex: many modes,
options, and negotiation steps. Two points matter for you. First, it is the
canonical ``heavyweight'' contrast that motivated WireGuard's minimalism.
Second, the IETF's interim post-quantum answer for IPsec, \textbf{RFC 8784},
mixes a \emph{pre-shared post-quantum key} into IKEv2 --- which is conceptually
identical to your \textbf{PSK-slot baseline}: get PQ confidentiality without
touching the asymmetric handshake, at the cost of out-of-band key distribution.

%------------------------------------------------------------------------------
\subsection{OpenVPN}
\seethesis{Background contrast; representative TLS-based userspace VPN.}
\nutshell{A long-established userspace VPN that rides on TLS/OpenSSL and can
tunnel at layer 2 or 3 --- flexible and battle-tested, but heavier and slower
than WireGuard.}

OpenVPN builds its tunnel on top of a TLS session (historically via OpenSSL),
typically over UDP. Its flexibility (certificates, pluggable ciphers, TAP layer-2
or TUN layer-3 modes) is also its weight: a larger codebase, more configuration,
and more overhead per packet than WireGuard. In your framing it is the
``general-purpose, TLS-based'' point of comparison, against which WireGuard's
fixed, lean Noise design stands out --- and it inherits TLS's post-quantum story
rather than defining its own.

%------------------------------------------------------------------------------
\subsection{TLS}
\seethesis{\S3.5 (related work: PQ via TLS PSK delivery); industry precedent for hybrids.}
\nutshell{The transport-layer security protocol of the web; TLS 1.3 already
ships a deployed hybrid X25519+ML-KEM-768 key exchange, which is the strongest
real-world precedent for your design.}

TLS secures byte streams \emph{above} the transport layer (it sits on TCP, or on
QUIC at the application layer). TLS 1.3 reduced the handshake to one round trip
and pruned legacy options. The reason it belongs in your defense is the
precedent: in 2024 Chrome and Cloudflare turned on a \textbf{hybrid} key
exchange named \texttt{X25519MLKEM768} --- the exact same primitives you combine,
combined the same way (concatenate the two shared secrets, feed them through the
KDF). So your thesis is not exotic; it transplants an approach the wider industry
has already validated at scale into a protocol (WireGuard) that did not yet have
it. Related-work \S3.5 also discusses delivering a PQ secret to WireGuard
\emph{through} a TLS channel.

\keyq{``Why not just rely on TLS being post-quantum?'' TLS protects different
traffic (application byte streams), not the WireGuard tunnel itself. A WireGuard
session's confidentiality depends on \emph{WireGuard's} handshake; fixing it
requires changing that handshake, which is the contribution.}

%------------------------------------------------------------------------------
\subsection{MTU: the IPv4/IPv6 maximum transmission unit}
\seethesis{\S5.5 (message size analysis) and \S9.2 (deployment).}
\nutshell{The largest packet a link will carry in one piece; post-quantum
handshake messages are bigger, so you have to show they still fit.}

The MTU is the largest payload a link-layer technology will carry in a single
frame --- classically 1500 bytes on Ethernet. Exceed it and the packet must be
fragmented or dropped. The two IP versions differ crucially:
\begin{itemize}
  \item \textbf{IPv4} routers may fragment in transit; the guaranteed minimum
  reassembly is small (576 bytes is the conventional safe figure).
  \item \textbf{IPv6} forbids in-network fragmentation --- only the sender may
  fragment --- and guarantees a \textbf{minimum MTU of 1280 bytes} on every
  link.
\end{itemize}
Your post-quantum messages grow because ML-KEM-768 adds an encapsulation key
($\approx$1184~B) and a ciphertext ($\approx$1088~B). The key result you defend:
the PQ initiation ($\approx$1332~B) and response ($\approx$1180~B) each still fit
inside a single 1500-byte Ethernet datagram, and even within the 1280-byte IPv6
floor the response fits and the initiation needs only minimal fragmentation.
Fragmentation only becomes a concern below a path MTU of roughly 1360 bytes.

%------------------------------------------------------------------------------
\subsection{PMTUD (Path MTU Discovery)}
\seethesis{\S9.2 (deployment considerations and failure modes).}
\nutshell{The mechanism hosts use to learn the smallest MTU along a path; it
fails silently (``blackholes'') when the ICMP signals it relies on are
filtered.}

PMTUD lets a sender discover the smallest MTU on the path to a destination.
In IPv4 the sender sets the ``Don't Fragment'' bit; a router that needs to
fragment instead drops the packet and returns an ICMP ``fragmentation needed''
message carrying the next-hop MTU. IPv6 does the analogous thing with ICMPv6
``Packet Too Big.'' The well-known failure mode --- a \textbf{PMTUD blackhole}
--- happens when firewalls filter those ICMP messages: the oversized packets
just vanish with no feedback, and the connection hangs. This matters to you
because if a path MTU were small enough to fragment the PQ handshake, and ICMP
were filtered, the handshake could stall. Your defense is that the messages are
sized to avoid fragmentation on normal paths in the first place, so PMTUD
pathologies do not bite. Note also: lowering the WireGuard interface MTU does
\emph{not} help, because the large packets are the UDP-encapsulated handshake on
the outer path, not inner tunnel traffic.

%==============================================================================
\section{Classical Cryptographic Primitives}
%==============================================================================
These are the pieces WireGuard already uses, plus the hard problems they rest on.
Know which ones quantum computers threaten (the public-key ones) and which they
do not (the symmetric ones) --- that split is the entire rationale for the thesis.

%------------------------------------------------------------------------------
\subsection{Diffie--Hellman key exchange}
\seethesis{\S2.1 (the quantum threat to public-key crypto); foundation for X25519.}
\nutshell{The 1976 method by which two parties derive a shared secret over a
public channel; WireGuard's whole handshake is a sequence of these.}

Diffie--Hellman (DH) lets two parties who have never met agree on a shared secret
visible to no eavesdropper. In the classic finite-field form, public parameters
are a prime $p$ and generator $g$; Alice sends $g^a \bmod p$, Bob sends
$g^b \bmod p$, and both compute $g^{ab}$. An eavesdropper sees $g^a$ and $g^b$
but recovering $g^{ab}$ requires solving the \textbf{discrete logarithm problem}
(DLP). DH provides \emph{key agreement}, not authentication --- on its own it is
vulnerable to a meddler-in-the-middle, which is why WireGuard authenticates by
mixing long-term static keys into the handshake. Crucially, DH (in both
finite-field and elliptic-curve form) is broken by Shor's algorithm; this is the
quantum problem your thesis addresses.

%------------------------------------------------------------------------------
\subsection{The integer factorization problem}
\seethesis{\S2.1 (why classical public-key crypto falls to quantum).}
\nutshell{Given a large $N = p\cdot q$, recover the primes $p$ and $q$ ---
believed hard classically, easy for a quantum computer.}

Factoring underpins RSA. The best classical algorithm, the General Number Field
Sieve, runs in \emph{sub-exponential} time --- slow enough that 2048-bit moduli
are safe against classical attackers. The catch is ``classical.'' Shor's
algorithm factors in polynomial time on a quantum computer, collapsing the
security margin. You cite factorization as one of the two classical assumptions
(with the discrete log) that quantum computing dissolves.

%------------------------------------------------------------------------------
\subsection{RSA}
\seethesis{\S2.1 (representative classical public-key scheme under quantum threat).}
\nutshell{The 1977 public-key cryptosystem (encryption and signatures) whose
security reduces to integer factorization --- and which Shor's algorithm breaks.}

RSA (Rivest--Shamir--Adleman) was the first practical public-key system. A public
key is $(N, e)$ with $N = pq$; the private key is the matching exponent $d$.
Security rests on factoring being hard: knowing $N$ should not reveal $p$ and
$q$. WireGuard does \emph{not} use RSA --- it uses elliptic-curve crypto --- but
RSA is the textbook example of a factoring-based scheme and the headline victim
of Shor in resource estimates (e.g.\ Gidney--Eker\aa, 2019, estimate RSA-2048 in
hours on $\sim$20 million noisy qubits). It frames the threat your work responds
to.

%------------------------------------------------------------------------------
\subsection{The elliptic-curve discrete logarithm problem (ECDLP)}
\seethesis{\S2.1 and \S2.2 (the hardness behind X25519).}
\nutshell{Given points $P$ and $Q = kP$ on an elliptic curve, recover the scalar
$k$ --- the hard problem that makes X25519's small keys secure (classically).}

On an elliptic curve, the group operation is point addition; ``exponentiation''
becomes scalar multiplication $kP$. The ECDLP is recovering $k$ from $P$ and
$Q=kP$. Its key property: the best classical attacks (Pollard's rho) cost about
$\sqrt{n}$ for group order $n$, i.e.\ fully exponential, with \emph{no}
sub-exponential shortcut like the number field sieve. That is why elliptic-curve
keys can be far smaller than RSA keys at equal classical security --- 256-bit
curves versus 3072-bit RSA. But ECDLP, like ordinary DLP and factoring, also
falls to Shor's algorithm; in fact a 256-bit ECDLP needs \emph{fewer} logical
qubits than RSA-2048, so elliptic-curve crypto is arguably the more urgent quantum
target.

%------------------------------------------------------------------------------
\subsection{Curve25519, X25519, and the 32-byte key}
\seethesis{\S2.2 (WireGuard primitives), \S5.3 (where ML-KEM joins it in the hybrid).}
\nutshell{Curve25519 is the elliptic curve; X25519 is the Diffie--Hellman
function over it; the public key and shared secret are each 32 bytes --- this is
the classical half of your hybrid.}

\textbf{Curve25519} (Bernstein, 2006) is a Montgomery-form elliptic curve over
the prime field $\mathbb{F}_p$ with $p = 2^{255}-19$, designed for speed and for
resistance to implementation pitfalls (it is ``safe'' by construction --- no
invalid-curve or twist problems if used correctly, and amenable to constant-time
code). \textbf{X25519} is the Diffie--Hellman function built on it (RFC 7748):
each side has a 32-byte private scalar (with specific bits ``clamped'') and a
32-byte public key (the curve point's $u$-coordinate); the shared secret is the
32-byte $u$-coordinate of the scalar product. Those ubiquitous \textbf{32-byte}
quantities --- WireGuard public keys, X25519 outputs --- are 256-bit values. In
your hybrid, the X25519 shared secret is one of the two inputs concatenated
before the KDF; ML-KEM-768's 32-byte shared secret is the other. X25519 is the
component a quantum attacker would eventually defeat via Shor; ML-KEM is what
keeps the session safe if that happens.

\keyq{``Why keep X25519 at all if it's quantum-broken?'' Two reasons: (1) ML-KEM
is young and lattice assumptions are less battle-tested than ECDLP, so a hybrid
stays secure if \emph{either} holds; (2) it preserves classical security and
WireGuard's existing authentication unchanged. Hybrid = belt and braces, which is
exactly NIST's and industry's recommended migration posture.}

%------------------------------------------------------------------------------
\subsection{ChaCha20-Poly1305}
\seethesis{\S2.2 (transport encryption), \S8.3 (throughput is unaffected by your change).}
\nutshell{WireGuard's authenticated-encryption cipher for the data channel: a
software-fast stream cipher (ChaCha20) plus a one-time MAC (Poly1305); your work
does not touch it.}

ChaCha20-Poly1305 (Bernstein; RFC 8439) is an \textbf{AEAD} (Authenticated
Encryption with Associated Data) construction: ChaCha20 is a 256-bit-key stream
cipher and Poly1305 is a fast one-time message authenticator; together they give
confidentiality \emph{and} integrity in one pass. WireGuard chose it over
AES-GCM because it is fast and \emph{constant-time in pure software}, with no
dependence on AES hardware acceleration --- which matters on devices like the
Raspberry Pi. The defense-relevant point: this is symmetric crypto, so it is
\emph{not} broken by Shor (only mildly weakened by Grover), and your thesis
leaves the entire transport path on ChaCha20-Poly1305 untouched. That is why your
measured throughput is statistically unchanged --- only the one-time handshake
cost grows.

%------------------------------------------------------------------------------
\subsection{AES-GCM}
\seethesis{Contrast to ChaCha20-Poly1305 (\S2.2); explains a WireGuard design choice.}
\nutshell{The other mainstream AEAD: AES in counter mode plus the GHASH
authenticator --- excellent with hardware AES, which is why WireGuard chose
ChaCha20-Poly1305 instead for hardware-agnostic speed.}

AES-GCM (NIST SP 800-38D) combines AES in counter mode with the GHASH universal
hash for authentication. It is the dominant AEAD in TLS and is extremely fast
\emph{where AES-NI hardware instructions exist}. Its weaknesses in software are
the flip side: without hardware acceleration, constant-time AES is harder and
slower, and GCM is unforgiving of nonce reuse. This contrast explains WireGuard's
preference for ChaCha20-Poly1305 and is worth knowing because your Raspberry Pi
results lean on the ``no AES hardware needed'' argument. Like all symmetric
ciphers it is quantum-resistant with adequate key length.

%------------------------------------------------------------------------------
\subsection{BLAKE2s}
\seethesis{\S2.2 (hashing/MAC), \S5.4 (the KDF that mixes your hybrid secret).}
\nutshell{WireGuard's hash function --- fast, modern, and used for hashing, the
keyed MAC, and the KDF; it is what stirs your concatenated secret into the chain.}

BLAKE2s is a cryptographic hash (a member of the BLAKE2 family derived from the
SHA-3 finalist BLAKE) optimized for 8--32-bit platforms, producing up to a
256-bit digest. WireGuard uses it pervasively: as the transcript hash that binds
the handshake, as a \emph{keyed} MAC (the \texttt{mac1}/\texttt{mac2} cookie
fields), and as the compression function inside its HMAC-based KDF. For your
thesis it is the hash underneath the key derivation in \S5.4: when you
concatenate the X25519 and ML-KEM secrets, it is a BLAKE2s-based HMAC chain that
absorbs them into the chaining key. Being a hash, it is quantum-resistant
(Grover only halves preimage security, which 256-bit output absorbs).

%------------------------------------------------------------------------------
\subsection{HMAC and HKDF key derivation}
\seethesis{\S2.5 (concatenate-then-KDF) and \S5.4 (key derivation details).}
\nutshell{HMAC is a hash-based message authentication code; HKDF is the
``extract-then-expand'' key-derivation function built from it --- the machinery
that turns your two raw shared secrets into uniform session keys.}

\textbf{HMAC} (RFC 2104) builds a MAC from any hash $H$ via
$\mathrm{HMAC}(k,m) = H\big((k\oplus opad)\mathbin\| H((k\oplus ipad)\mathbin\| m)\big)$.
\textbf{HKDF} (Krawczyk; RFC 5869) is a KDF with two stages: \emph{extract}
condenses raw, possibly-non-uniform input keying material into a uniform
pseudorandom key, then \emph{expand} stretches that into as many output keys as
needed. WireGuard uses an HKDF-style HMAC-BLAKE2s chain throughout its handshake.

This is the heart of why ``concatenate-then-KDF'' is sound. You set the KDF input
keying material to the concatenation $ss_{\text{X25519}} \,\|\, ss_{\text{ML-KEM}}$.
A standard KDF security property (and the NIST SP 800-56C Rev.\,2 guidance for
combining secrets) guarantees the output is pseudorandom as long as \emph{at
least one} input secret is unknown to the attacker. That single property is the
formal backbone of your hybrid's ``secure if either primitive holds'' claim ---
no new cryptography invented, just a correct application of an established
combiner.

\keyq{``Why concatenate rather than XOR or nest the secrets?'' Concatenation
feeding a KDF is the construction NIST SP 800-56C analyses and the TLS hybrid
uses; it provably preserves security if either secret is strong. XOR can be
fragile if one secret is adversarially controlled; concatenation-then-KDF is the
conservative, standardized choice.}

%==============================================================================
\section{The Noise Protocol Framework and the WireGuard Handshake}
%==============================================================================

%------------------------------------------------------------------------------
\subsection{The Noise Protocol Framework}
\seethesis{\S2.3 (Noise) and \S5.3 (where you insert ML-KEM into the pattern).}
\nutshell{A toolkit (Trevor Perrin) for building handshake protocols out of a
fixed alphabet of Diffie--Hellman and mixing operations; WireGuard is one
specific Noise pattern.}

Noise is not a protocol but a \emph{framework} for constructing two-party
handshakes. A handshake is written as a short \emph{pattern} of tokens --- e.g.\
\texttt{e} (send an ephemeral key), \texttt{s} (send a static key), \texttt{ee},
\texttt{es}, \texttt{se}, \texttt{ss} (perform the corresponding
Diffie--Hellman) --- and the framework prescribes how each operation mixes into
two running values: a \emph{chaining key} (feeds the KDF) and a \emph{handshake
hash} (binds the transcript for authentication). The elegance is that security
properties (forward secrecy, identity hiding, authentication) follow from the
pattern. The catch for post-quantum work: Noise's vocabulary assumes
\emph{Diffie--Hellman}, and a KEM is not a DH (see \S4.1). So inserting ML-KEM
means stepping slightly outside vanilla Noise --- exactly the careful integration
your implementation performs and documents.

%------------------------------------------------------------------------------
\subsection{The Noise IK pattern (and WireGuard's IKpsk2)}
\seethesis{\S2.3, \S5.3 (hybrid handshake), \S5.2 / \S2.2.1 (the PSK slot = the ``psk2'').}
\nutshell{IK is the specific Noise handshake WireGuard uses: the initiator knows
the responder's static key in advance (K) and transmits its own static key,
encrypted, during the handshake (I).}

In Noise's two-letter naming, the first letter describes the \emph{initiator's}
static key and the second the \emph{responder's}:
\begin{itemize}
  \item \textbf{I} --- the initiator's static key is transmitted (\emph{immediately},
  encrypted) to the responder during the handshake.
  \item \textbf{K} --- the responder's static key is \emph{known} to the
  initiator beforehand (out of band; in WireGuard, from the peer's configured
  public key).
\end{itemize}
This fits a VPN perfectly: a client already has the server's public key in its
config (K), and proves its own identity in the first message (I). WireGuard
extends IK with an optional pre-shared key mixed at the second position, making it
effectively \textbf{IKpsk2}; that ``psk2'' slot is precisely the \textbf{PSK-slot
baseline} you compare against (Option~2). Your \textbf{hybrid} (Option~1) instead
threads ML-KEM into the existing two-message flow: the initiator carries the
ML-KEM encapsulation key, the responder returns the ciphertext, and the resulting
KEM shared secret is mixed into the same chaining key alongside the X25519 DH
results --- no third message added.

\keyq{``Did you add a round trip?'' No. ML-KEM's encapsulation key rides in the
existing initiation and its ciphertext rides in the existing response, so the
handshake stays 1-RTT. Only the message \emph{sizes} and the per-handshake
\emph{compute} grow; the message \emph{count} does not.}

%------------------------------------------------------------------------------
\subsection{Unknown-Key-Share (UKS) attacks}
\seethesis{\S7.4 (authentication and remaining limitations).}
\nutshell{An attack where two honest parties end up sharing a key but disagree
about \emph{who} they are talking to; prevented by binding both identities into
the transcript --- which WireGuard already does.}

In a UKS attack the adversary does not learn the key, but it tricks Alice into
believing she shares a key with the attacker while she actually shares it with
Bob (or vice versa) --- a confusion of identities rather than of key material.
The standard defense is to fold both parties' identities (their static public
keys) into the authenticated handshake transcript, so the derived key is
cryptographically committed to ``who talked to whom.'' WireGuard does this: static
public keys are mixed into the handshake hash. You raise UKS in the security
analysis to show you understand the authentication model and that your change
preserves it --- your modification adds confidentiality material (the KEM secret)
without weakening the identity binding. It also sets up the honest limitation:
\emph{authentication} still rests on X25519 static keys, so it is not itself
post-quantum (addressed next, and in \S7.4/\S9.3).

%==============================================================================
\section{The Quantum Threat}
%==============================================================================

%------------------------------------------------------------------------------
\subsection{Shor's algorithm}
\seethesis{\S2.1 (the quantum threat) --- the reason the thesis exists.}
\nutshell{The 1994 quantum algorithm that solves factoring and discrete logs
(including elliptic-curve) in polynomial time, breaking RSA, DH, and X25519.}

Peter Shor's algorithm (1994) is the single result that endangers today's
public-key crypto. Its engine is quantum \emph{period-finding}: factoring $N$
reduces to finding the period of the function $x \mapsto a^x \bmod N$, and the
\textbf{Quantum Fourier Transform} extracts that period efficiently --- turning a
sub-exponential classical problem into a polynomial-time quantum one. The same
period-finding machinery solves the discrete logarithm in any finite abelian
group, which sweeps in finite-field DH \emph{and} the elliptic-curve discrete log
behind X25519. So a sufficiently large, fault-tolerant quantum computer would
break WireGuard's key exchange outright.

The essential caveat: ``sufficiently large and fault-tolerant.'' Current devices
are noisy and far too small; published estimates put RSA-2048 at millions of
physical qubits with error correction. No such machine exists today --- which is
exactly why the threat model is not ``tomorrow'' but \emph{harvest-now,
decrypt-later} (\S4.3).

\keyq{``If no quantum computer can do this yet, why migrate now?'' Because
encrypted traffic captured today can be stored and decrypted later once the
machine exists. Data with multi-year confidentiality needs is already at risk;
the migration must precede the computer, not follow it. NIST/NSA timelines target
the early 2030s.}

%------------------------------------------------------------------------------
\subsection{What Shor breaks --- and what it does not (Grover)}
\seethesis{\S2.1; justifies why only the handshake, not the transport, needs PQ.}
\nutshell{Shor demolishes public-key (asymmetric) crypto; symmetric crypto and
hashes are only mildly dented by Grover's algorithm, so doubling key sizes
suffices there.}

This split is the load-bearing argument for your whole scope:
\begin{itemize}
  \item \textbf{Asymmetric / public-key} (RSA, DH, ECDH/X25519, ECDSA) ---
  \emph{broken} by Shor. These need replacing. WireGuard's key exchange lives
  here.
  \item \textbf{Symmetric \& hashes} (ChaCha20, AES, BLAKE2s, SHA-2) --- only
  \emph{weakened} by Grover's algorithm, which gives at most a quadratic speedup
  for brute-force search. Effective security halves: a 256-bit cipher retains
  $\sim$128-bit post-quantum security, which is ample.
\end{itemize}
Therefore the only quantum-vulnerable part of WireGuard is the \textbf{handshake}
(the X25519 key agreement); the ChaCha20-Poly1305 transport is fine as-is. This
is precisely why your contribution targets the handshake \emph{only} and leaves
the data plane untouched --- a scoping decision you can now justify from first
principles.

%------------------------------------------------------------------------------
\subsection{Harvest-now, decrypt-later (HNDL)}
\seethesis{\S1 (introduction / motivation) and \S2.1.}
\nutshell{The real, present-day threat: an adversary records encrypted traffic
now and decrypts it once a quantum computer arrives --- so long-lived secrets
need protection before the machine exists.}

HNDL (also ``store-now-decrypt-later'') reframes the quantum threat from a future
event into a present obligation. An adversary with storage and patience captures
today's encrypted sessions and simply waits; when a cryptographically-relevant
quantum computer arrives, the recorded X25519 handshakes are broken and the
session keys --- hence all that captured traffic --- are recovered. Anything that
must stay confidential for years (state secrets, health and financial records,
long-term credentials) is therefore \emph{already} exposed. This is the urgency
argument and it is specifically a \emph{confidentiality} threat, which is why a
hybrid that protects the key exchange is valuable even though it leaves classical
\emph{authentication} in place: forging authentication requires a quantum
computer at handshake time and cannot be done retroactively against recorded
traffic.

%==============================================================================
\section{Post-Quantum Cryptography: Standards and Constructions}
%==============================================================================
The core of your contribution's ``new half.'' Be fluent here above all.

%------------------------------------------------------------------------------
\subsection{NIST}
\seethesis{\S2.4 (ML-KEM is a NIST/FIPS standard).}
\nutshell{The U.S.\ National Institute of Standards and Technology --- the body
that ran the post-quantum competition and publishes the FIPS standards your work
implements.}

NIST is a U.S.\ federal agency that, among much else, standardizes cryptography
through its Federal Information Processing Standards (FIPS) publications --- AES
(FIPS 197) and SHA-3 (FIPS 202) came through NIST processes, and so did ML-KEM
(FIPS 203). Its imprimatur is why ML-KEM, rather than raw Kyber, is the thing to
deploy: it is the vetted, frozen, interoperable specification. When the committee
hears ``FIPS 203,'' they hear ``the standard,'' and your choice to target it
rather than a research draft signals engineering maturity.

%------------------------------------------------------------------------------
\subsection{The NIST PQC Standardization Project}
\seethesis{\S2.4 and \S3 (related work situates competing schemes).}
\nutshell{The multi-year open competition (2016--) that evaluated post-quantum
algorithms and selected the first standards in 2022, published as FIPS in 2024.}

NIST launched the Post-Quantum Cryptography Standardization Project in 2016,
soliciting public submissions for quantum-resistant key-establishment and
signature schemes. The timeline you should be able to sketch:
\begin{itemize}
  \item \textbf{2017} --- 69 candidates accepted into Round 1.
  \item \textbf{2019 / 2020} --- narrowed to Round 2 (26) then Round 3 (7
  finalists + 8 alternates).
  \item \textbf{July 2022} --- first selections: \textbf{CRYSTALS-Kyber} for
  key encapsulation, plus CRYSTALS-Dilithium, Falcon, and SPHINCS+ for
  signatures.
  \item \textbf{August 2024} --- standards published: \textbf{FIPS 203}
  (ML-KEM, from Kyber), FIPS 204 (ML-DSA, from Dilithium), FIPS 205 (SLH-DSA,
  from SPHINCS+).
  \item A separate ``Round 4'' continued evaluating additional KEMs for
  diversity (BIKE, Classic McEliece, HQC; SIKE was eliminated after being
  broken), with HQC selected in 2025.
\end{itemize}
The process matters rhetorically: ML-KEM is not a lab curiosity but the survivor
of years of public cryptanalysis --- which is the assurance your hybrid leans on
for its post-quantum half.

%------------------------------------------------------------------------------
\subsection{The Module Learning With Errors (MLWE) problem}
\seethesis{\S2.4 (the hardness assumption under ML-KEM).}
\nutshell{The lattice problem ML-KEM's security rests on: distinguishing
noisy linear equations over a module from random --- believed hard for classical
\emph{and} quantum computers.}

Learning With Errors (LWE) asks you to recover a secret $\mathbf{s}$ from many
samples $(\mathbf{a}_i, \langle \mathbf{a}_i, \mathbf{s}\rangle + e_i)$ where each
$e_i$ is small random ``noise'' --- equivalently, to distinguish those noisy
inner products from uniform randomness. The noise is what makes Gaussian
elimination fail and the problem hard; its hardness connects to worst-case
lattice problems. Plain LWE has large keys, so practical schemes use structure:
\textbf{Ring-LWE} works over a polynomial ring (compact but more structure
assumed), and \textbf{Module-LWE} interpolates between the two by working over a
rank-$k$ \emph{module} of that ring. Kyber/ML-KEM uses MLWE over
$R_q = \mathbb{Z}_q[X]/(X^{256}+1)$ with module rank $k$ tied to the security
level. The selling point: MLWE lets you tune security by changing $k$ while
reusing the same ring arithmetic, and no efficient quantum algorithm is known
against it --- unlike the abelian-group problems Shor demolishes.

%------------------------------------------------------------------------------
\subsection{CRYSTALS-Kyber, ML-KEM, and Kyber-768}
\seethesis{\S2.4 and \S5.3 (the post-quantum half of the hybrid).}
\nutshell{Kyber is the MLWE-based KEM that won the NIST competition; ML-KEM is
its standardized form (FIPS 203); ML-KEM-768 / ``Kyber-768'' is the security
level you use.}

\textbf{CRYSTALS-Kyber} (Cryptographic Suite for Algebraic Lattices) is a
key-encapsulation mechanism built on MLWE. \textbf{ML-KEM} is its NIST-blessed
final form; the names are often used interchangeably, but ``ML-KEM'' denotes the
exact FIPS-203 algorithm (with a few tweaks from the competition Kyber). As a KEM
it exposes three operations:
\begin{itemize}
  \item \texttt{KeyGen} $\to$ (encapsulation key / public, decapsulation key /
  secret);
  \item \texttt{Encaps}(encapsulation key) $\to$ (ciphertext, shared secret) ---
  run by the party that has the peer's public key;
  \item \texttt{Decaps}(decapsulation key, ciphertext) $\to$ shared secret.
\end{itemize}
\textbf{ML-KEM-768} (``Kyber-768'') is the middle parameter set: module rank
$k=3$, targeting NIST security \textbf{category~3} (roughly AES-192). Its sizes
--- which you should have at your fingertips because they drive the MTU analysis
--- are an encapsulation key of \textbf{1184 bytes}, a ciphertext of
\textbf{1088 bytes}, and a \textbf{32-byte} shared secret. You chose the 768
level as the balanced, widely-recommended default (it is the level in TLS's
\texttt{X25519MLKEM768}), strong enough for category~3 without the larger
messages of ML-KEM-1024.

\keyq{``Why ML-KEM-768 and not 512 or 1024?'' 768 is NIST category~3 and the
de-facto industry default (it is the level chosen for the deployed TLS hybrid).
512 trades security margin for size; 1024 buys category~5 at noticeably larger
messages. 768 is the balanced choice and keeps your handshake within one Ethernet
datagram.}

%------------------------------------------------------------------------------
\subsection{FIPS 203}
\seethesis{\S2.4 (the standard your implementation targets).}
\nutshell{The August 2024 NIST standard that defines ML-KEM precisely --- the
specification your RustCrypto \texttt{ml-kem} dependency implements.}

FIPS 203, ``Module-Lattice-Based Key-Encapsulation Mechanism Standard,''
finalized on 13 August 2024, is the authoritative definition of ML-KEM: the
algorithms, the three parameter sets (512/768/1024), the exact byte encodings,
and the internal hashing. Targeting a published FIPS --- rather than a moving
research target --- is what makes your implementation interoperable and your
choice defensible. The crate you depend on (\texttt{ml-kem}, RustCrypto) tracks
this standard, which is why correctness can be checked against NIST's official
test vectors (\S6 below).

%------------------------------------------------------------------------------
\subsection{The Number-Theoretic Transform (NTT)}
\seethesis{\S2.4; explains why ML-KEM is fast and underlies your latency numbers.}
\nutshell{An FFT over a finite field that turns slow polynomial multiplication
into cheap coefficient-wise products --- the reason Kyber is fast, and the reason
its parameters $q=3329$, $n=256$ look the way they do.}

Kyber's arithmetic lives in $R_q = \mathbb{Z}_q[X]/(X^{256}+1)$, and multiplying
two degree-255 polynomials naively is $O(n^2)$. The NTT is the finite-field
analogue of the FFT: it evaluates a polynomial at the powers of a root of unity
so that multiplication becomes \emph{pointwise}, giving $O(n\log n)$. Kyber's
parameters are chosen \emph{precisely} to make this work: $q = 3329$ is prime
with $q \equiv 1 \pmod{256}$, so a length-256 transform exists. (A subtlety worth
knowing: because $q \not\equiv 1 \pmod{512}$, $X^{256}+1$ does not split into
linear factors, so Kyber uses an ``incomplete'' NTT down to 128 quadratic factors
with base-case products --- but the speed win stands.) Public keys even travel in
NTT domain. This is the mechanistic reason your ML-KEM operations are fast (tens
of microseconds on the M4); your benchmark chapter measures the practical payoff.

%------------------------------------------------------------------------------
\subsection{IND-CPA and IND-CCA2 security}
\seethesis{\S2.4 (KEM security) and \S7.2 (hybrid security argument).}
\nutshell{The two standard ``indistinguishability'' security definitions; ML-KEM
must meet the stronger one (IND-CCA2), because a real attacker can submit chosen
ciphertexts.}

These are game-based security \emph{definitions}, framed as an attacker trying to
distinguish encryptions:
\begin{itemize}
  \item \textbf{IND-CPA} (chosen-plaintext): the attacker may encrypt plaintexts
  of its choosing but cannot ask for decryptions. The minimal sensible bar.
  \item \textbf{IND-CCA2} (adaptive chosen-ciphertext): the attacker additionally
  has a \emph{decryption oracle} and may use it adaptively (except on the
  challenge). This models the real world, where an adversary can send crafted
  ciphertexts to a peer and observe behavior.
\end{itemize}
A KEM used in a live protocol must be \textbf{IND-CCA2}, because a network
attacker can inject ciphertexts. ML-KEM's underlying public-key scheme (K-PKE) is
only IND-CPA; the upgrade to IND-CCA2 is done by the Fujisaki--Okamoto transform
(next). In your security chapter, the KEM's IND-CCA2 guarantee is one of the two
pillars; the KDF combiner then carries that strength into the hybrid session key.

%------------------------------------------------------------------------------
\subsection{The Fujisaki--Okamoto (FO) transformation}
\seethesis{\S2.4 (how ML-KEM achieves CCA security).}
\nutshell{The generic recipe that upgrades a merely IND-CPA encryption scheme
into an IND-CCA2 KEM by re-encrypting on decapsulation and rejecting anything
that doesn't match --- this is how ML-KEM gets its strong security.}

The FO transform turns a weak (IND-CPA) public-key scheme into a strong
(IND-CCA2) KEM. The idea: the shared secret is derived deterministically from the
message, so on \texttt{Decaps} the receiver \emph{re-encrypts} the recovered
message and checks it reproduces the received ciphertext exactly. A mismatch
means tampering, and the KEM refuses to output the real key. ML-KEM uses the
modern variant with \textbf{implicit rejection}: instead of returning a visible
error (which can leak information), it returns a \emph{pseudorandom} key derived
from a secret seed, so a chosen-ciphertext attacker learns nothing from
malformed queries. This re-encryption check is also why decapsulation costs a bit
more than encapsulation --- relevant when you read your per-operation benchmarks.

\keyq{``Where does CCA security come from in ML-KEM?'' From the FO transform
applied to the IND-CPA scheme K-PKE, with implicit rejection. The decapsulator
re-encrypts and compares; on failure it returns a pseudorandom key, leaking
nothing.}

%------------------------------------------------------------------------------
\subsection{The road not taken: Classic McEliece and Saber}
\seethesis{\S2.4 / \S3 (context for why ML-KEM, not an alternative).}
\nutshell{Two other NIST KEM candidates --- code-based McEliece (ultra-conservative
but enormous keys) and lattice-based Saber (close runner-up to Kyber) --- useful
as foils for ``why ML-KEM.''}

\textbf{Classic McEliece} is a \emph{code-based} KEM (security from
error-correcting codes, an assumption studied since 1978 and untouched by
quantum attacks). Its virtue is conservatism; its problem is a public key of
hundreds of kilobytes to a megabyte --- impossible to carry in a 1500-byte
WireGuard handshake. \textbf{Saber} is a lattice KEM very close in spirit to
Kyber but based on Module Learning-With-\emph{Rounding} with power-of-two moduli
(so it avoids the NTT); it was a Round-3 finalist that NIST did not select. Both
are worth a sentence because they show you understand the design space: McEliece's
giant keys make it a non-starter for an MTU-bound handshake, and Saber lost to
Kyber on NIST's balance of security, performance, and analysis --- so ML-KEM-768
is the natural, standards-aligned, size-feasible choice for WireGuard.

%==============================================================================
\section{Assurance: Verification and Testing}
%==============================================================================
How anyone trusts that the cryptography is correct --- the standard's, the
library's, and yours.

%------------------------------------------------------------------------------
\subsection{EasyCrypt and formal verification}
\seethesis{\S2.4 / \S7 (why ML-KEM as a primitive can be trusted).}
\nutshell{A proof assistant for machine-checking cryptographic security proofs;
ML-KEM has been formally verified with it, which is part of why the primitive is
trustworthy and not your burden to re-prove.}

EasyCrypt is an interactive theorem prover specialized for the \emph{game-based}
proofs cryptographers write by hand --- it mechanically checks the sequence of
``game hops'' that bounds an attacker's advantage, removing the human errors that
have historically crept into paper proofs. It matters to your defense as a
boundary marker: the security of ML-KEM \emph{as a primitive} has been the subject
of formal-verification efforts (e.g.\ machine-checked Kyber/ML-KEM
implementations in the Jasmin/Formosa-crypto line of work using EasyCrypt), so you
inherit a strongly-vetted building block. Your contribution is the \emph{correct
integration} of that block into WireGuard, not a re-proof of the lattice
mathematics --- a clean scoping statement if asked ``did you prove ML-KEM
secure?''

%------------------------------------------------------------------------------
\subsection{NIST Known Answer Tests (KATs)}
\seethesis{\S6.5 (test harness) --- how you verify your integration is correct.}
\nutshell{Deterministic test vectors --- fixed seed in, exact expected bytes out
--- that confirm a crypto implementation matches the standard.}

A Known Answer Test pins down a deterministic computation: given a fixed seed (so
the otherwise-random key generation and encapsulation become reproducible), the
implementation must produce exactly the published keys, ciphertext, and shared
secret. NIST ships KAT files with FIPS 203 for precisely this. Their value is
that they catch byte-level mistakes --- wrong encoding, endianness, parameter
mix-ups --- that ordinary ``it ran'' testing misses. In your test harness, KAT
agreement (via the \texttt{ml-kem} crate's conformance) gives confidence that the
ML-KEM operations you wired into the handshake are the \emph{real} ML-KEM, and
your own integration tests then confirm two patched peers actually complete a
hybrid handshake and derive matching keys.

%------------------------------------------------------------------------------
\subsection{Historical security exploits in cryptographic code}
\seethesis{\S6 (rationale for a Rust implementation and a minimal, gated diff).}
\nutshell{The catalogue of real-world cryptographic failures --- memory bugs and
logic slips --- that motivates implementing in memory-safe Rust and keeping the
change small and optional.}

A short rogues' gallery you can cite to justify engineering choices:
\begin{itemize}
  \item \textbf{Heartbleed} (OpenSSL, 2014) --- a missing bounds check let
  attackers read adjacent memory, leaking private keys. A \emph{memory-safety}
  failure in C.
  \item \textbf{Apple ``goto fail''} (2014) --- a duplicated \texttt{goto}
  skipped TLS signature verification entirely. A \emph{logic} failure.
  \item \textbf{Debian OpenSSL PRNG} (2008) --- a well-meaning edit gutted
  entropy, making keys guessable. A \emph{review/process} failure.
\end{itemize}
The lessons map directly onto your decisions: BoringTun is \textbf{Rust}, whose
memory safety eliminates the Heartbleed class by construction; you keep the
post-quantum change \textbf{small, auditable, and behind a \texttt{pq} feature
flag} so the default build is byte-for-byte the upstream code (no risk to existing
users); and you lean on \textbf{standardized primitives with KATs and formal
verification} rather than hand-rolling crypto, which is where the logic failures
live. This is your answer to ``why should we trust your code?''

%==============================================================================
\section{Implementation Concepts}
%==============================================================================

%------------------------------------------------------------------------------
\subsection{C FFI (Foreign Function Interface)}
\seethesis{\S6.1--\S6.2 (BoringTun architecture and technology stack).}
\nutshell{The mechanism by which Rust exposes a C-callable API; BoringTun ships
one, and your \texttt{pq} flag is designed to leave that ABI untouched in the
default build.}

An FFI lets code in one language call code compiled from another; the C ABI is
the lingua franca. Rust exposes C-callable functions with \texttt{extern "C"} and
\texttt{\#[no\_mangle]}, and such boundaries are \texttt{unsafe} because the
compiler cannot verify the foreign side's assumptions (lifetimes, null pointers,
thread safety). BoringTun is consumed as a library through exactly such a C FFI
(it is embedded into other applications, not only run as a daemon). Why this
matters to you: keeping the post-quantum work behind a compile-time feature means
the \emph{default} build presents the same FFI surface and the same wire format
as upstream --- existing integrators are unaffected, and only an explicitly
PQ-enabled build changes behavior. That is a deployment-safety argument the
committee will appreciate.

%------------------------------------------------------------------------------
\subsection{Rust \texttt{no\_std}}
\seethesis{\S6.2 (technology stack); relevant to constrained-device deployability.}
\nutshell{A Rust mode that drops the OS-dependent standard library and uses only
the core language, enabling crypto to run on embedded targets --- which is why a
KEM library can reach routers and small devices.}

By default Rust programs link \texttt{std}, which assumes an operating system
(heap allocator, files, threads). A crate marked \texttt{\#![no\_std]} uses only
\texttt{core} (the allocation-free language essentials), optionally adding
\texttt{alloc} for heap types --- letting it run on microcontrollers and inside
kernels. RustCrypto primitives, including the \texttt{ml-kem} family, are written
to be \texttt{no\_std}-capable. The relevance to your thesis is forward-looking:
although BoringTun itself targets hosted platforms (it uses \texttt{std}), the
fact that the post-quantum primitive is \texttt{no\_std}-friendly means the
approach is portable down to constrained devices --- consistent with your
Raspberry Pi evaluation and with WireGuard's ambition to run everywhere.

%==============================================================================
\section{Evaluation Methodology}
%==============================================================================
How your performance claims are made trustworthy rather than anecdotal.

%------------------------------------------------------------------------------
\subsection{Criterion benchmarking}
\seethesis{\S6.5 (test/bench harness) and all of \S8 (results).}
\nutshell{Rust's statistical benchmarking library --- it warms up, samples many
iterations, detects outliers, and reports a mean with a confidence interval
rather than a single noisy number.}

Criterion is the de-facto statistics-driven benchmarking framework for Rust. Its
methodology is the point: it runs a \emph{warm-up} phase (to settle caches and
CPU frequency), then collects many timed samples, fits and reports a \emph{mean
estimate with a confidence interval}, detects \emph{outliers}, and can compare
against a saved baseline to flag regressions. This is what lifts your numbers from
``I ran it once and got 247\,$\mu$s'' to a defensible statistical statement. You
extended BoringTun's existing Criterion-based tests to measure the handshake under
the three configurations (vanilla, PSK-slot, hybrid) and the individual ML-KEM
operations, on both the Apple M4 and the Raspberry Pi~5.

%------------------------------------------------------------------------------
\subsection{Confidence intervals}
\seethesis{\S8.1 (latency results are reported as mean $\pm$ 95\% CI).}
\nutshell{A range that quantifies measurement uncertainty: a 95\% CI means the
procedure captures the true mean 95\% of the time --- so non-overlapping
intervals signal a real difference, not noise.}

A confidence interval expresses how sure you are about an estimate. A
\textbf{95\% CI} for the mean is constructed so that, across repeated
experiments, 95\% of such intervals would contain the true mean. Criterion derives
these by \emph{bootstrap resampling} the collected timings (no assumption of a
normal distribution). For your defense the practical reading is comparative: when
the hybrid's latency interval sits clearly above vanilla's and the two do not
overlap, the $\sim$64\% increase on the M4 (and $\sim$25\% on the Pi) is a
\emph{statistically real} effect, not measurement jitter --- whereas the
near-identical vanilla and PSK-slot intervals overlap, correctly showing the PSK
approach adds no asymmetric-handshake cost. Reporting the interval, not just the
mean, is what makes the comparison honest.

\keyq{``Are your performance differences significant?'' Yes --- they are reported
as means with 95\% bootstrap confidence intervals from Criterion over many
samples; the hybrid's interval is disjoint from vanilla's, so the overhead is a
real effect, while vanilla vs.\ PSK-slot overlap, correctly indicating no
handshake-asymmetric cost.}

\vspace{1.2em}
\hrule height 1pt
\vspace{0.6em}
\noindent\textbf{\color{deepcyan}Closing orientation.} If you remember nothing
else: \emph{Shor breaks the public-key handshake but not the symmetric transport};
\emph{a KEM is not a Diffie--Hellman}, so ML-KEM is woven into Noise IK as
encapsulation-key-out / ciphertext-back rather than a DH exchange; \emph{the two
secrets are concatenated and run through the BLAKE2s KDF}, which is secure if
\emph{either} holds; and \emph{the cost is a bigger, slower handshake but an
unchanged data plane}, measured honestly with Criterion confidence intervals.
Everything in this companion is a more detailed footnote to those four facts.

\end{document}
